\scanpage{19}
\deleted{\unsure{$E$}}

$A$ courbe elliptique

$Y=A\times P^{\unsure{\deleted{1}|1}}$\qquad $P$ courbe non. s. projective conn.

$\Theta_0=\Delta-A\times(0)-(0)A$ \deleted{diviseur \unsure{sur}} diviseur corr. sur $A\times A$ définissant l'autodualité

$\Theta=\Theta_0\times P_{\unsure{p}}^{\unsure{\deleted{1}|1}}$\qquad diviseur correspondant sur $A\times A\times P^{\unsure{\deleted{1}|1}}=A\times Y$

$c$, classe de diviseur sur $Y$, défini pr un \deleted{fibré} faisceau inv. $\underline L$ comme \inserted{$c=\operatorname{cl}(L)$}

$X=\overline{V(\underline L)}$ (Clôture projective).
\diagram{\begin{tikzcd}[ampersand replacement=\&]
Y^2\arrow[r,"i"] \& X^3
\end{tikzcd}}
section nulle.

le cycle $\Theta$ sur $A\times Y$ définit
\[
z=(\operatorname{id}_A\times i)_*(\Theta)
\]
cycle de codim $2$ sur \deleted{$X$} $A\times X^3$. On veut calculer, si $s:A\times X^3\mathbin{-}A$ est la projection, $s_*(z^2)$, et on va donner (pour des choix convenables de $c$) des cas où il est $>0$, d'autres où il est $<0$.

Pour ceci, notons que si $i'=\operatorname{id}_A\times i$, on a
\[
z^2=i'_*(\Theta)^2=i'_*(\Theta^2c')
\]
où
\[
\begin{aligned}
c'&=i'^*(i'_*(1))=c^1\ (\text{faisceau normal pr }i')\\
&=1_A\times c,
\end{aligned}
\]
$c$ étant le faisceau normal pr \deleted{\unsure{$Y$}} $i:Y\mathbin{-}X$.

\margin{\diagram{\begin{tikzcd}[ampersand replacement=\&]
A\times Y\arrow[dr,"t"'] \& \& A\times X\arrow[ll,"i'"']\arrow[dl,"s"]\\
\& A \&
\end{tikzcd}}}

\mspage{18}
